\chapter{Complexity — Frag}
%%% generated by the blueprint pipeline from TCSlib.Complexity.SpaceComplexity.Machines.Frag
\section{Overview}

Declarations of the Lean module \texttt{TCSlib.Complexity.SpaceComplexity.Machines.Frag}.

\section{Declarations}

\begin{definition}[Half-step register action]
\lean{Complexity.LogProg.halfLAct}
\label{Complexity.LogProg.halfLAct}
\leanok
\uses{Complexity.LogProg.regAct, Turing.Action}
Given a register, a carry symbol that may be absent, states $h_F,h_T,$ and $next$, and
the symbol $rd$ read from the register (or blank), the action moves the register head
left and writes the carry symbol when $rd$ is present, entering $h_T$ if it is true and
$h_F$ if it is false. When $rd$ is blank, it leaves the register unchanged, moves its
head right, and enters $next$.
\end{definition}

\begin{definition}[Half-step state selector]
\lean{Complexity.LogProg.halfSt}
\label{Complexity.LogProg.halfSt}
\leanok
For three states $h_0,h_F,h_T$, define a map from $\text{Option}(\text{Bool})$ to states
by sending `none` to $h_0$, `some false` to $h_F$, and `some true` to $h_T$.
\end{definition}

\begin{lemma}[Register halving run]
\lean{Complexity.LogProg.halfL_run}
\label{Complexity.LogProg.halfL_run}
\leanok
\uses{Complexity.LogProg.apply_regAct, Complexity.LogProg.bufferTape_erase_last, Complexity.LogProg.halfLAct, Complexity.LogProg.halfSt, Complexity.LogProg.regCfg, Complexity.LogProg.regCfg_read, Complexity.LogProg.rrun, Complexity.LogProg.rrun_add, Complexity.LogProg.rrun_one_reg, Complexity.LogProg.rrun_succ_left, Complexity.LogProg.update_bufferTape_cons, Turing.Cfg, Turing.FinTM.bufferTape}
\difficulty{6}
Let $P$ be a register program with oracle and input, let $r$ be a register, and let
$h_R,h_0,h_F,h_T,$ and $next$ be states. Assume that at $h_R$ every transition is the
specified end-marker action, at $h_0,h_F,h_T$ every transition is the specified halving
action for the scanned symbol, and none of these four states is a call state. For any
configuration $c$, Boolean word $w$, and integer $k$ with $k\le |w|$, initialize the
program in state $h_0$, $h_F$, or $h_T$ according as the bit of $w$ at index $k$ is
absent, $0$, or $1$; place on register $r$ the word obtained by deleting that bit, with
its head at $k-1$. After $k+1$ steps the configuration is in state $next$, register $r$
contains the tail of $w$ with its head at $0$, and all other configuration data agree
with $c$. At each time $t<k+1$, the configuration has some state and contents for
register $r$ and head position $q$ such that the state is not a call state and $-1\le
q\le k-1$.
\end{lemma}

\begin{lemma}[Register halving computation]
\lean{Complexity.LogProg.half_run}
\label{Complexity.LogProg.half_run}
\leanok
\uses{Complexity.LogProg.bits_half, Complexity.LogProg.halfL_run, Complexity.LogProg.regCfg, Complexity.LogProg.rrun, Complexity.LogProg.rrun_add, Complexity.LogProg.toEnd_run, Turing.Cfg, Turing.FinTM.bufferTape}
\difficulty{4}
Let $P$ be a register program with designated register $r$ and states
$h_R,h_0,h_F,h_T,next$. Suppose its transitions at these states implement the stated
end-marker and leftward halving actions for each possible scanned symbol, and none of
these four states is a call state. For every configuration $c$ and natural number $n$,
there is a time $T$ at which the program, started in state $h_R$ with register $r$
containing the binary word for $n$ and its head at zero, reaches state $next$ with that
register containing the binary word for $\lfloor n/2\rfloor$ and its head at zero. At
every earlier time $t<T$, its configuration differs from $c$ only in the state, tape,
and head position of register $r$; the state is not a call state and the head position
lies between $-1$ and the length of the binary word for $n$, inclusive.
\end{lemma}

\begin{definition}[Configuration with Two Head Positions Set]
\lean{Complexity.LogProg.eqCfg}
\label{Complexity.LogProg.eqCfg}
\leanok
\uses{Turing.Cfg}
Let $c$ be a configuration of an $m$-tape machine on input $x$, let $s$ be a state, let
$r_1,r_2$ be work-tape indices, and let $k$ be an integer. The resulting configuration
has state $s$, the same input-head position, work-tape contents, and output as $c$, and
has both work-tape heads at indices $r_1$ and $r_2$ positioned at $k$; if the indices
coincide, that head is positioned at $k$.
\end{definition}

\begin{definition}[Two-Tape Movement Action]
\lean{Complexity.LogProg.mv2Act}
\label{Complexity.LogProg.mv2Act}
\leanok
\uses{Turing.Action}
Let $r_1,r_2$ be work-tape indices, let $v\in\{-1,0,1\}$, and let $s$ be a state. The
action leaves the input head stationary, writes nothing and moves each of the two
indexed work-tape heads by $v$, leaves every other work-tape head stationary without
writing, emits no output symbol, and enters state $s$.
\end{definition}

\begin{lemma}[Simultaneous movement of two registers]
\lean{Complexity.LogProg.apply_mv2Act}
\label{Complexity.LogProg.apply_mv2Act}
\leanok
\uses{Complexity.LogProg.eqCfg, Complexity.LogProg.mv2Act, Turing.Action.apply, Turing.Cfg, Turing.moveInputPos_zero}
\difficulty{4}
Let $c$ be a configuration with $m$ Boolean work tapes, state set $\Lambda$, and fixed
input $x$. For states $s,s'\in\Lambda$, register indices $r_1,r_2\in\{0,\ldots,m-1\}$,
integer $k$, and movement $mv\in\{-1,0,1\}$, applying the action that moves both
designated register heads by $mv$ and enters state $s'$ to the configuration obtained
from $c$ by setting the state to $s$ and both designated head positions to $k$ yields
the configuration with state $s'$, unchanged input position, tapes, and output, and both
designated head positions set to $k+mv$.
\end{lemma}

\begin{definition}[Comparison-Dependent Tape Action]
\lean{Complexity.LogProg.eqCAct}
\label{Complexity.LogProg.eqCAct}
\leanok
\uses{Complexity.LogProg.mv2Act, Turing.Action}
Let $r_1,r_2$ be work-tape indices, let $e_C,e_{By},e_{Bn}$ be states, and let $a,b$ be
either Boolean symbols or blank. If $a=b$ and both are blank, the action moves the two
indexed work-tape heads one cell left and enters $e_{By}$. If $a=b$ and they are not
blank, it moves them one cell right and enters $e_C$; if $a\ne b$, it moves them one
cell left and enters $e_{Bn}$. In every case it leaves the input head stationary, writes
nothing, and emits no output symbol.
\end{definition}

\begin{definition}[Boolean-Dependent Tape Action]
\lean{Complexity.LogProg.eqBAct}
\label{Complexity.LogProg.eqBAct}
\leanok
\uses{Complexity.LogProg.mv2Act, Turing.Action}
Let $r_1,r_2$ be work-tape indices, let $e_B,l$ be states, and let $a$ be either a
Boolean symbol or blank. If $a$ is a symbol, the action moves the two indexed work-tape
heads one cell left and enters $e_B$; if $a$ is blank, it moves them one cell right and
enters $l$. In both cases it leaves the input head stationary, writes nothing, and emits
no output symbol.
\end{definition}

\begin{lemma}[Scanned symbol after head repositioning]
\lean{Complexity.LogProg.eqCfg_read₁}
\leanok
\uses{Complexity.LogProg.eqCfg, Turing.Cfg, Turing.Cfg.workTapeSymbols}
\difficulty{1}
Let $r_1$ and $r_2$ be distinct work-tape registers. For any configuration $c$, state
$s$, and integer $k$, form the configuration with state $s$, the same tapes and input
head as $c$, and both register-head positions set to $k$. The symbol then scanned on
register $r_1$ is the content of register $r_1$ in $c$ at cell $k$.
\end{lemma}

\begin{lemma}[Reading the Second Register]
\lean{Complexity.LogProg.eqCfg_read₂}
\leanok
\uses{Complexity.LogProg.eqCfg, Turing.Cfg, Turing.Cfg.workTapeSymbols}
\difficulty{1}
For a configuration $c$, a state $s$, distinct work-tape indices $r_1,r_2$, and an
integer $k$, change the state to $s$ and place both work-tape heads at $k$, leaving the
tapes and all other configuration data unchanged. The symbol read at the second head is
then the content of tape $r_2$ at cell $k$.
\end{lemma}

\begin{lemma}[One ordinary transition from an encoded configuration]
\lean{Complexity.LogProg.rrun_one_eq}
\label{Complexity.LogProg.rrun_one_eq}
\leanok
\uses{Complexity.LogProg.eqCfg, Complexity.LogProg.rrun, Complexity.LogProg.rrun_one, Complexity.LogProg.rstep_noncall, Turing.Action.apply, Turing.Cfg, Turing.Cfg.inputSymbol, Turing.Cfg.workTapeSymbols, Turing.MultiTapeTM.step}
\difficulty{3}
Let $P$ be a register program with oracle, let $r_1,r_2$ be register indices, and let
$c$ be a configuration with state set $\Lambda$. For a state $s\in\Lambda$ and position
$k\in\mathbb Z$, form the configuration with state $s$ and both register heads at
position $k$. If $P$ has no oracle call at $s$, then one step of its run from this
configuration equals applying the machine transition prescribed by $P$ for state $s$,
the scanned input symbol, and the scanned work-tape symbols.
\end{lemma}

\begin{lemma}[Run comparing two register words]
\lean{Complexity.LogProg.eqC_run}
\label{Complexity.LogProg.eqC_run}
\leanok
\uses{Complexity.LogProg.apply_mv2Act, Complexity.LogProg.eqCAct, Complexity.LogProg.eqCfg, Complexity.LogProg.eqCfg_read₁, Complexity.LogProg.eqCfg_read₂, Complexity.LogProg.rrun, Complexity.LogProg.rrun_add, Complexity.LogProg.rrun_one_eq, Turing.Cfg, Turing.FinTM.bufferTape}
\difficulty{6}
Let $P$ be a register program with oracle, and let $r_1,r_2$ be distinct registers.
Suppose the transitions at states $e_C,e_{By},e_{Bn}$ implement, respectively, the
bitwise comparison of the two registers, the comparison continuation leading to yes
state $yes$, and the comparison continuation leading to no state $no$; assume none of
these three states is a call state. If the two registers initially store words
$w_1,w_2$, then for any words $u_1,u_2,pre$ with $w_1=pre\,u_1$ and $w_2=pre\,u_2$,
there are a time $T$ and an integer position $p$ such that $|pre|-1\le p<|w_1|$, and the
run from the comparison state with both heads at position $|pre|$ reaches state $e_{By}$
if $w_1=w_2$ and state $e_{Bn}$ otherwise, with both heads at position $p$. Before time
$T$, the run remains in state $e_C$, with both heads at a common position $q$ satisfying
$|pre|\le q\le |w_1|$.
\end{lemma}

\begin{lemma}[Run through a binary register]
\lean{Complexity.LogProg.eqB_run}
\label{Complexity.LogProg.eqB_run}
\leanok
\uses{Complexity.LogProg.apply_mv2Act, Complexity.LogProg.eqBAct, Complexity.LogProg.eqCfg, Complexity.LogProg.eqCfg_read₁, Complexity.LogProg.rrun, Complexity.LogProg.rrun_add, Complexity.LogProg.rrun_one_eq, Complexity.LogProg.rrun_succ_left, Turing.Cfg, Turing.FinTM.bufferTape}
\difficulty{5}
Let $P$ be a register program with oracle, and let $r_1,r_2$ be distinct registers.
Suppose the transition at state $e_B$ performs the binary comparison step for $r_1$,
leading to $l$ when the scanned bit is blank and otherwise moving left while remaining
at $e_B$, and suppose $e_B$ is not a call state. If the contents of register $r_1$
initially store the word $w_1$, then for every natural number $n$ and integer $p$ with
$p+1=n$ and $p<|w_1|$, running for $n+1$ steps from state $e_B$ with both register heads
at position $p$ reaches state $l$ with both heads at position zero. At every earlier
time $t<n+1$, the configuration has state $e_B$ and both heads at some integer position
$q$ with $-1\le q\le p$.
\end{lemma}

\begin{lemma}[Execution of the Equality Program]
\lean{Complexity.LogProg.eq_run}
\label{Complexity.LogProg.eq_run}
\leanok
\uses{Complexity.LogProg.bits_injective, Complexity.LogProg.eqB_run, Complexity.LogProg.eqC_run, Complexity.LogProg.eqCfg, Complexity.LogProg.rrun, Complexity.LogProg.rrun_add, Turing.Cfg, Turing.FinTM.bufferTape}
\difficulty{5}
Let $P$ be a register program with oracle, let $r_1$ and $r_2$ be distinct registers,
and let $e_C,e_{B_y},e_{B_n},y,n$ be states. Assume the transition function at these
three states satisfies the specified equality-check and branching action rules,
respectively, for every input symbol and tuple of scanned work-tape symbols, and assume
the program makes no oracle call at any of these states. If registers $r_1,r_2$
initially store the buffer tapes for the binary words of natural numbers $a,b$, then
there is a time $T$ at which execution from the configuration with state $e_C$ and both
heads at cell $0$ reaches the same configuration with state $y$ when $a=b$ and state $n$
otherwise. At every earlier time $t<T$, the configuration has some state $s$, both heads
at a common cell $q$, no oracle call is made at $s$, and $-1\le q\le \ell$, where $\ell$
is the length of the binary word for $a$.
\end{lemma}
